\documentclass[a4paper, noanswers, addpoints, 12pt]{exam}
\usepackage[T1]{fontenc}
\usepackage[utf8]{inputenc}
\usepackage{lmodern}
\usepackage[french]{babel}
\usepackage{xcolor}
\usepackage{listings}
\usepackage{graphicx}
\usepackage{array}
\usepackage{amsmath}
\usepackage{amssymb}

% Configuration de la coloration syntaxique (style du document joint)
\definecolor{codegreen}{rgb}{0,0.6,0}
\definecolor{codegray}{rgb}{0.5,0.5,0.5}
\definecolor{codepurple}{rgb}{0.58,0,0.82}
\definecolor{backcolour}{rgb}{0.95,0.95,0.92}
 
\lstdefinestyle{mystyle}{
    backgroundcolor=\color{backcolour},   
    commentstyle=\color{codegreen},
    keywordstyle=\color{magenta},
    numberstyle=\tiny\color{codegray},
    stringstyle=\color{codepurple},
    basicstyle=\footnotesize\ttfamily,
    breakatwhitespace=false,         
    breaklines=true,                 
    captionpos=b,                    
    keepspaces=true,                 
    numbers=left,                    
    numbersep=5pt,                  
    showspaces=false,                
    showstringspaces=false,
    showtabs=false,                  
    tabsize=2
}
\lstset{style=mystyle, language=Python}

\title{\vspace{-2cm}Examen Term NSI - Tri, recherche et complexité}

\usepackage[a4paper,top=2.3cm,bottom=2.8cm,left=2cm,right=2cm,headsep=0.6cm,footskip=1.2cm]{geometry}
\begin{document}
\runningfooter{}{\footnotesize Babash - NSI Tani Malandi et Johanne Bakalara \textperiodcentered{} CC BY-SA 4.0 \textperiodcentered{} page \thepage}{}
\firstpagefooter{}{\footnotesize Babash - NSI Tani Malandi et Johanne Bakalara \textperiodcentered{} CC BY-SA 4.0 \textperiodcentered{} page \thepage}{}

\maketitle
\thispagestyle{headandfoot}

\pagestyle{headandfoot}
\thispagestyle{headandfoot}
\firstpageheader{\makebox[\textwidth]{Nom : \dotfill\ Pr\'enom : \dotfill\ Classe : \makebox[3cm]{\dotfill}}}{}{}
\vspace{-3cm}


\section*{Consignes:}
\begin{center} 
\fbox{\parbox{\linewidth}{
\begin{center}
La clarté de la rédaction et les explications font partie intégrante de la notation. Pour les questions demandant du code, respectez la syntaxe Python et soyez attentifs aux détails.
\end{center}


}}
\end{center}

\vsword{Note}
\hsword{Note}

\section*{Barème :}
\begin{center}
  \par\bigskip
  \gradetable[h]
\end{center}

\vspace{0.5cm}

\begin{questions}

% --- COMPLEXITÉ THÉORIQUE ---
\question \textbf{Complexité et acceptabilité}
\begin{parts}
    \part[2] Classez par ordre croissant d'efficacité : $O(n^2)$, $O(1)$, $O(n \log n)$, $O(n)$, $O(2^n)$, $O(\log n)$.
    \fillwithdottedlines{1cm}
    \part[1] Un algorithme de tri fusion ($O(n \log n)$) met 1ms pour trier 1000 éléments. Si on passe à $10^6$ éléments, le temps devient-il "infini" (plusieurs jours) ou reste-t-il de l'ordre de quelques secondes ? Justifiez intuitivement.
    \fillwithdottedlines{1.5cm}
\end{parts}

% --- QUESTION RELIER (CONSERVÉE) ---
\question[2\half] \textbf{Principes de fonctionnement} 
Reliez chaque algorithme à la description qui correspond à son principe fondamental :
\begin{center}
\begin{tabular}{r c l}
Recherche linéaire $\bullet$ & & $\bullet$ Diviser pour régner : trier deux moitiés puis les réunir. \\
Recherche dichotomique $\bullet$ & & $\bullet$ Insérer la valeur actuelle à sa place dans la partie triée. \\
Tri par insertion $\bullet$ & & $\bullet$ Parcourir un par un tous les éléments. \\
Tri par sélection $\bullet$ & & $\bullet$ Diviser l'espace de recherche par deux à chaque étape. \\
Tri fusion $\bullet$ & & $\bullet$ Chercher le plus petit élément pour le placer au début. \\
\end{tabular}
\end{center}

\question[2\half] Donner la classe de complexité de chacun des algorithmes suivants \\
\begin{itemize}
    \item Recherche linéaire: $O$(\dotfill) Recherche dichotomique: $O$(\dotfill)
    \item Tri par insertion: $O$(\dotfill) Tri par sélection: $O$(\dotfill) Tri fusion: $O$(\dotfill)
\end{itemize}

\newpage

% --- QUESTION CODES (CONSERVÉE & OPTIMISÉE) ---
\question[2\half] \textbf{Identification de code} \\

\begin{minipage}[t]{0.48\textwidth}
\begin{lstlisting}[basicstyle=\scriptsize\ttfamily, title={Code A}]
while g <= d:
    m = (g + d) // 2
    if t[m] == v:
        return m
    elif t[m] < v:
        g = m + 1
    else:
        d = m - 1
\end{lstlisting}
\begin{lstlisting}[basicstyle=\scriptsize\ttfamily, title={Code B}]
for i in range(1, len(t)):
    k = t[i]
    j = i - 1
    while j >= 0 and t[j] > k:
        t[j+1] = t[j]
        j -= 1
    t[j+1] = k
\end{lstlisting}
\end{minipage}
\hfill
\begin{minipage}[t]{0.48\textwidth}
\begin{lstlisting}[basicstyle=\scriptsize\ttfamily, title={Code C}]
for i in range(len(t)):
    if t[i] == v:
        return i
\end{lstlisting}
\begin{lstlisting}[basicstyle=\scriptsize\ttfamily, title={Code D}]
for i in range(len(t)-1):
    m = i
    for j in range(i+1, len(t)):
        if t[j] < t[m]: m = j
    t[i], t[m] = t[m], t[i]
\end{lstlisting}
\begin{lstlisting}[basicstyle=\scriptsize\ttfamily, title={Code E}]
if len(t) <= 1:
    return t
m = len(t) // 2
g = tri(t[:m])
d = tri(t[m:])
return fusion(g, d)
\end{lstlisting}
\end{minipage}

Attribuez à chaque bloc de code Python l'algorithme correspondant :
\textit{Recherche linéaire, Recherche dichotomique, Tri par insertion, Tri par sélection et Tri fusion}
\begin{itemize}
    \item A:\dotfill B:\dotfill
    \item C:\dotfill D:\dotfill
    \item E:\dotfill
\end{itemize}

% --- ANALYSE DE CODE & COMPLEXITÉ (MUTUALISATION) ---
\question \textbf{Analyse de la recherche}
\begin{parts}
    \part[1] Quelle condition obligatoire le tableau \texttt{t} doit-il respecter pour que la \textbf{Recherche dichotomique} fonctionne ? \fillwithdottedlines{0.8cm}
    \part[1] On exécute un \textbf{algorithme de recherche dichotomique} avec \texttt{t = [2, 5, 8, 12, 18]} et \texttt{v = 12}. Donnez les valeurs successives de la variable \texttt{m}. \\
    \fillwithdottedlines{0.8cm}
    \part[2] Si le tableau contient $n = 1\ 000\ 000$ éléments, quel est le nombre maximum de comparaisons pour : \\
    \textit{Aide :} $log_2(250\ 000)\approx18$\\
    - La \textbf{Recherche linéaire} (Pire cas) : \dotfill \\
    - La \textbf{Recherche dichotomique} (Pire cas, environ) : \dotfill
\end{parts}

\newpage

% --- PARTIE 6 : FOCUS TRI PAR INSERTION (6.5 points) ---

\question \textbf{Étude approfondie du Tri par Insertion} \\
On s'appuie sur le \textbf{Code B} vu précédemment pour répondre aux questions suivantes.

\begin{parts}
    \part[1\half] \textbf{Déroulé de l'algorithme} : Soit le tableau \texttt{t = [7, 2, 5, 3]}. 
    Donnez l'état du tableau après chaque itération de la boucle externe \texttt{for i} :
    \begin{itemize}
        \item État initial : \texttt{[7, 2, 5, 3]}
        \item Fin itération $i=1$ : \dotfill
        \item Fin itération $i=2$ : \dotfill
        \item Fin itération $i=3$ : \dotfill
    \end{itemize}

    \part[1] \textbf{Invariant de boucle} : Parmi les propositions suivantes, laquelle est l'invariant de la boucle \texttt{for} ? (Cochez la bonne réponse)
    \begin{checkboxes}
        \choice À l'étape $i$, les $i$ derniers éléments sont les plus grands du tableau.
        \choice À l'étape $i$, la tranche \texttt{t[0...i]} contient les mêmes éléments qu'au départ, mais triés.
        \choice À l'étape $i$, l'élément \texttt{t[i]} est à sa place définitive dans le tableau.
    \end{checkboxes}

    \part\textbf{Analyse des performances} : 
    \begin{subparts}
        \subpart[1] Quel est le \textbf{meilleur des cas} pour ce tri (quel type de liste) ? Quelle est alors sa complexité ?
        \fillwithdottedlines{1cm}
        
        \subpart[1] Quel est le \textbf{pire des cas} pour ce tri ? Quelle est alors sa complexité ?
        \fillwithdottedlines{1cm}
        
        \subpart[1] Pourquoi dit-on que cet algorithme est "efficace" sur une liste presque triée par rapport au tri par sélection ? 
        \fillwithdottedlines{1.2cm}
    \end{subparts}
\end{parts}
\end{questions}

\end{document}